% ============================================================================
%  Introduction -- % ============================================================================
%  Introduction -- % ============================================================================
%  Introduction -- % ============================================================================
%  Introduction -- \input{sections/introduction}
%  Contribution bullets that depend on results carry \todo notes: keep, rephrase or drop
%  them once Tables in sections/results.tex are filled.
% ============================================================================
\section{Introduction}
\label{sec:intro}

Feed is one of the largest operating costs in aquaculture, and how it is dispensed directly affects
productivity, water quality and fish welfare~\cite{atoum2014feeding,yang2021smartfish}. Most feeding
machines still run on fixed schedules or on the judgement of an operator, so fish are underfed when
appetite is high and overfed when it is low. The uneaten feed is wasted and degrades the water.
\emph{Fish feeding intensity assessment} (FFIA) estimates the appetite of a group of fish during feeding,
usually as an ordinal level from \textsc{None} to \textsc{Strong}. It is the sensing step needed to close
the loop between feeding machines and the fish~\cite{an2021behaviour,zhou2019ffia}.

Most FFIA systems are vision-based. They classify single images~\cite{zhou2019ffia,yang2021effnetb2,zhang2023msif}
or exploit motion through optical flow and 3-D convolutions~\cite{ubina2021ffia}. Vision captures
aggregation and surface agitation directly, but it suffers from low light, reflections, turbidity
and occlusion, and video models are expensive to run on site. Acoustic FFIA uses the sounds of feeding,
such as chewing, tail slapping and splashing, picked up by a hydrophone~\cite{cui2022affia3k}. It is
unaffected by lighting and cheap to process, but it is sensitive to the bubble and pump noise that is
ever-present in recirculating systems. The two modalities fail under different conditions, which
makes them natural partners. The release of AV-FFIA~\cite{cui2024uffia}, 27,000 synchronised
audio-visual clips of \textit{Oplegnathus punctatus} labelled into four intensities, made this
combination measurable. U-FFIA showed that fusion beats either modality, especially in noise, and
that a unified model can lower the cost~\cite{cui2024uffia}.

Accuracy on AV-FFIA has since risen quickly, with later works reporting 95--98\%
\cite{acafnet2025,dynagate2026,hmmfnet2026,hmkd2026,ldmf2026}. Three problems remain open, however.
\emph{First, efficiency.} The fusion models of~\cite{cui2024uffia} cost 15--69\,GMACs and 21--48\,M
parameters per two-second clip. Recent compact models help, but parameter count, memory and the
cost of transmitting models to edge devices on farms still matter.
\emph{Second, how and where to fuse.} Fusion is usually done once, by concatenation, cross-attention
or attention bottlenecks, at a fixed depth. On AV-FFIA these choices behave very differently: the
attention bottleneck (MBT~\cite{nagrani2021mbt}) reached only 72.4\% against 91.7\% for cross-attention
in~\cite{cui2024uffia}, and there has been no systematic study of the depth at which audio and video
should interact.
\emph{Third, evaluation.} Published numbers rest on incompatible protocols: the official random split,
5-fold cross-validation, and filtered subsets of the data. Moreover, every split is drawn at clip level,
although the clips are consecutive two-second segments of one-minute recordings. Near-duplicate
neighbours can therefore appear in both training and test, which may inflate accuracy.

Hypercomplex neural networks offer a principled answer to the first two problems.
A layer over a hypercomplex algebra (complex numbers, quaternions, octonions, sedenions) shares one set
of weight blocks across all components through the algebra's multiplication table. This divides the
parameter count by the algebra's dimension~$n$ while still mixing every component with every other
\cite{parcollet2019qrnn,zhang2021phm,grassucci2022phnn}. In spatiotemporal tasks such networks match
or exceed their real-valued counterparts at a fraction of the parameters
\cite{bojesomo2024hypercomplex,bojesomo2023octonion,bojesomo2020sedenion}. The same property suggests a
new view of fusion. If the audio view, the video view and their interaction are taken as the
\emph{components} of one hypercomplex number, a single Hamilton product mixes all of them through
shared weights. Cross-modal interaction then becomes a property of the algebra rather than something
a dense layer must learn from limited data.

We build on this view and propose \textbf{HyperAVNet}, a hypercomplex audio-visual network for FFIA. Its
contributions are as follows.
\begin{itemize}
  \item A \textbf{fully hypercomplex audio-visual architecture}: Hamilton-product weight sharing in the
  convolutional backbones, attention, fusion and heads. It has 2.14\,M parameters, $3.7\times$ fewer than
  the identical real-valued network, at 1.18\,GMACs per clip; that is $10\times$ fewer parameters and
  $13\times$ fewer MACs than U-FFIA.
  \item \textbf{Hypercomplex component fusion (HCF)}: modality-specific views and their multiplicative
  agreement are encoded as the components of a quaternion (or octonion) and mixed by one Hamilton
  product. Its parameter cost does not depend on the number of interaction terms.
  \item A \textbf{deep-to-early fusion study with zero-initialised cross-modal feedback}, which shows where
  audio-visual interaction helps on AV-FFIA. \todo{one-line finding, e.g. ``fusing from the second stage
  onwards is best''}
  \item \textbf{Multi-level distillation} at the logit, embedding and relation levels from pretrained
  audio and video teachers, together with a \textbf{teacher ablation} over ten backbone pairs that
  tests whether a stronger teacher yields a stronger student. \todo{one-line finding}
  \item A \textbf{reproducible evaluation} under three protocols (the official split, 5-fold
  cross-validation, and a new \emph{session-disjoint} split), with a noise and missing-modality
  robustness study. \todo{State the headline: accuracy relative to the strongest comparable result,
  and the gap between clip-level and session-disjoint accuracy.}
\end{itemize}

The rest of the paper is organised as follows. Section~\ref{sec:related} reviews FFIA, audio-visual
fusion, hypercomplex networks and efficient models. Section~\ref{sec:method} presents HyperAVNet,
Section~\ref{sec:setup} the experimental set-up, and Section~\ref{sec:results} the results and
ablations.

%  Contribution bullets that depend on results carry \todo notes: keep, rephrase or drop
%  them once Tables in sections/results.tex are filled.
% ============================================================================
\section{Introduction}
\label{sec:intro}

Feed is one of the largest operating costs in aquaculture, and how it is dispensed directly affects
productivity, water quality and fish welfare~\cite{atoum2014feeding,yang2021smartfish}. Most feeding
machines still run on fixed schedules or on the judgement of an operator, so fish are underfed when
appetite is high and overfed when it is low. The uneaten feed is wasted and degrades the water.
\emph{Fish feeding intensity assessment} (FFIA) estimates the appetite of a group of fish during feeding,
usually as an ordinal level from \textsc{None} to \textsc{Strong}. It is the sensing step needed to close
the loop between feeding machines and the fish~\cite{an2021behaviour,zhou2019ffia}.

Most FFIA systems are vision-based. They classify single images~\cite{zhou2019ffia,yang2021effnetb2,zhang2023msif}
or exploit motion through optical flow and 3-D convolutions~\cite{ubina2021ffia}. Vision captures
aggregation and surface agitation directly, but it suffers from low light, reflections, turbidity
and occlusion, and video models are expensive to run on site. Acoustic FFIA uses the sounds of feeding,
such as chewing, tail slapping and splashing, picked up by a hydrophone~\cite{cui2022affia3k}. It is
unaffected by lighting and cheap to process, but it is sensitive to the bubble and pump noise that is
ever-present in recirculating systems. The two modalities fail under different conditions, which
makes them natural partners. The release of AV-FFIA~\cite{cui2024uffia}, 27,000 synchronised
audio-visual clips of \textit{Oplegnathus punctatus} labelled into four intensities, made this
combination measurable. U-FFIA showed that fusion beats either modality, especially in noise, and
that a unified model can lower the cost~\cite{cui2024uffia}.

Accuracy on AV-FFIA has since risen quickly, with later works reporting 95--98\%
\cite{acafnet2025,dynagate2026,hmmfnet2026,hmkd2026,ldmf2026}. Three problems remain open, however.
\emph{First, efficiency.} The fusion models of~\cite{cui2024uffia} cost 15--69\,GMACs and 21--48\,M
parameters per two-second clip. Recent compact models help, but parameter count, memory and the
cost of transmitting models to edge devices on farms still matter.
\emph{Second, how and where to fuse.} Fusion is usually done once, by concatenation, cross-attention
or attention bottlenecks, at a fixed depth. On AV-FFIA these choices behave very differently: the
attention bottleneck (MBT~\cite{nagrani2021mbt}) reached only 72.4\% against 91.7\% for cross-attention
in~\cite{cui2024uffia}, and there has been no systematic study of the depth at which audio and video
should interact.
\emph{Third, evaluation.} Published numbers rest on incompatible protocols: the official random split,
5-fold cross-validation, and filtered subsets of the data. Moreover, every split is drawn at clip level,
although the clips are consecutive two-second segments of one-minute recordings. Near-duplicate
neighbours can therefore appear in both training and test, which may inflate accuracy.

Hypercomplex neural networks offer a principled answer to the first two problems.
A layer over a hypercomplex algebra (complex numbers, quaternions, octonions, sedenions) shares one set
of weight blocks across all components through the algebra's multiplication table. This divides the
parameter count by the algebra's dimension~$n$ while still mixing every component with every other
\cite{parcollet2019qrnn,zhang2021phm,grassucci2022phnn}. In spatiotemporal tasks such networks match
or exceed their real-valued counterparts at a fraction of the parameters
\cite{bojesomo2024hypercomplex,bojesomo2023octonion,bojesomo2020sedenion}. The same property suggests a
new view of fusion. If the audio view, the video view and their interaction are taken as the
\emph{components} of one hypercomplex number, a single Hamilton product mixes all of them through
shared weights. Cross-modal interaction then becomes a property of the algebra rather than something
a dense layer must learn from limited data.

We build on this view and propose \textbf{HyperAVNet}, a hypercomplex audio-visual network for FFIA. Its
contributions are as follows.
\begin{itemize}
  \item A \textbf{fully hypercomplex audio-visual architecture}: Hamilton-product weight sharing in the
  convolutional backbones, attention, fusion and heads. It has 2.14\,M parameters, $3.7\times$ fewer than
  the identical real-valued network, at 1.18\,GMACs per clip; that is $10\times$ fewer parameters and
  $13\times$ fewer MACs than U-FFIA.
  \item \textbf{Hypercomplex component fusion (HCF)}: modality-specific views and their multiplicative
  agreement are encoded as the components of a quaternion (or octonion) and mixed by one Hamilton
  product. Its parameter cost does not depend on the number of interaction terms.
  \item A \textbf{deep-to-early fusion study with zero-initialised cross-modal feedback}, which shows where
  audio-visual interaction helps on AV-FFIA. \todo{one-line finding, e.g. ``fusing from the second stage
  onwards is best''}
  \item \textbf{Multi-level distillation} at the logit, embedding and relation levels from pretrained
  audio and video teachers, together with a \textbf{teacher ablation} over ten backbone pairs that
  tests whether a stronger teacher yields a stronger student. \todo{one-line finding}
  \item A \textbf{reproducible evaluation} under three protocols (the official split, 5-fold
  cross-validation, and a new \emph{session-disjoint} split), with a noise and missing-modality
  robustness study. \todo{State the headline: accuracy relative to the strongest comparable result,
  and the gap between clip-level and session-disjoint accuracy.}
\end{itemize}

The rest of the paper is organised as follows. Section~\ref{sec:related} reviews FFIA, audio-visual
fusion, hypercomplex networks and efficient models. Section~\ref{sec:method} presents HyperAVNet,
Section~\ref{sec:setup} the experimental set-up, and Section~\ref{sec:results} the results and
ablations.

%  Contribution bullets that depend on results carry \todo notes: keep, rephrase or drop
%  them once Tables in sections/results.tex are filled.
% ============================================================================
\section{Introduction}
\label{sec:intro}

Feed is one of the largest operating costs in aquaculture, and how it is dispensed directly affects
productivity, water quality and fish welfare~\cite{atoum2014feeding,yang2021smartfish}. Most feeding
machines still run on fixed schedules or on the judgement of an operator, so fish are underfed when
appetite is high and overfed when it is low. The uneaten feed is wasted and degrades the water.
\emph{Fish feeding intensity assessment} (FFIA) estimates the appetite of a group of fish during feeding,
usually as an ordinal level from \textsc{None} to \textsc{Strong}. It is the sensing step needed to close
the loop between feeding machines and the fish~\cite{an2021behaviour,zhou2019ffia}.

Most FFIA systems are vision-based. They classify single images~\cite{zhou2019ffia,yang2021effnetb2,zhang2023msif}
or exploit motion through optical flow and 3-D convolutions~\cite{ubina2021ffia}. Vision captures
aggregation and surface agitation directly, but it suffers from low light, reflections, turbidity
and occlusion, and video models are expensive to run on site. Acoustic FFIA uses the sounds of feeding,
such as chewing, tail slapping and splashing, picked up by a hydrophone~\cite{cui2022affia3k}. It is
unaffected by lighting and cheap to process, but it is sensitive to the bubble and pump noise that is
ever-present in recirculating systems. The two modalities fail under different conditions, which
makes them natural partners. The release of AV-FFIA~\cite{cui2024uffia}, 27,000 synchronised
audio-visual clips of \textit{Oplegnathus punctatus} labelled into four intensities, made this
combination measurable. U-FFIA showed that fusion beats either modality, especially in noise, and
that a unified model can lower the cost~\cite{cui2024uffia}.

Accuracy on AV-FFIA has since risen quickly, with later works reporting 95--98\%
\cite{acafnet2025,dynagate2026,hmmfnet2026,hmkd2026,ldmf2026}. Three problems remain open, however.
\emph{First, efficiency.} The fusion models of~\cite{cui2024uffia} cost 15--69\,GMACs and 21--48\,M
parameters per two-second clip. Recent compact models help, but parameter count, memory and the
cost of transmitting models to edge devices on farms still matter.
\emph{Second, how and where to fuse.} Fusion is usually done once, by concatenation, cross-attention
or attention bottlenecks, at a fixed depth. On AV-FFIA these choices behave very differently: the
attention bottleneck (MBT~\cite{nagrani2021mbt}) reached only 72.4\% against 91.7\% for cross-attention
in~\cite{cui2024uffia}, and there has been no systematic study of the depth at which audio and video
should interact.
\emph{Third, evaluation.} Published numbers rest on incompatible protocols: the official random split,
5-fold cross-validation, and filtered subsets of the data. Moreover, every split is drawn at clip level,
although the clips are consecutive two-second segments of one-minute recordings. Near-duplicate
neighbours can therefore appear in both training and test, which may inflate accuracy.

Hypercomplex neural networks offer a principled answer to the first two problems.
A layer over a hypercomplex algebra (complex numbers, quaternions, octonions, sedenions) shares one set
of weight blocks across all components through the algebra's multiplication table. This divides the
parameter count by the algebra's dimension~$n$ while still mixing every component with every other
\cite{parcollet2019qrnn,zhang2021phm,grassucci2022phnn}. In spatiotemporal tasks such networks match
or exceed their real-valued counterparts at a fraction of the parameters
\cite{bojesomo2024hypercomplex,bojesomo2023octonion,bojesomo2020sedenion}. The same property suggests a
new view of fusion. If the audio view, the video view and their interaction are taken as the
\emph{components} of one hypercomplex number, a single Hamilton product mixes all of them through
shared weights. Cross-modal interaction then becomes a property of the algebra rather than something
a dense layer must learn from limited data.

We build on this view and propose \textbf{HyperAVNet}, a hypercomplex audio-visual network for FFIA. Its
contributions are as follows.
\begin{itemize}
  \item A \textbf{fully hypercomplex audio-visual architecture}: Hamilton-product weight sharing in the
  convolutional backbones, attention, fusion and heads. It has 2.14\,M parameters, $3.7\times$ fewer than
  the identical real-valued network, at 1.18\,GMACs per clip; that is $10\times$ fewer parameters and
  $13\times$ fewer MACs than U-FFIA.
  \item \textbf{Hypercomplex component fusion (HCF)}: modality-specific views and their multiplicative
  agreement are encoded as the components of a quaternion (or octonion) and mixed by one Hamilton
  product. Its parameter cost does not depend on the number of interaction terms.
  \item A \textbf{deep-to-early fusion study with zero-initialised cross-modal feedback}, which shows where
  audio-visual interaction helps on AV-FFIA. \todo{one-line finding, e.g. ``fusing from the second stage
  onwards is best''}
  \item \textbf{Multi-level distillation} at the logit, embedding and relation levels from pretrained
  audio and video teachers, together with a \textbf{teacher ablation} over ten backbone pairs that
  tests whether a stronger teacher yields a stronger student. \todo{one-line finding}
  \item A \textbf{reproducible evaluation} under three protocols (the official split, 5-fold
  cross-validation, and a new \emph{session-disjoint} split), with a noise and missing-modality
  robustness study. \todo{State the headline: accuracy relative to the strongest comparable result,
  and the gap between clip-level and session-disjoint accuracy.}
\end{itemize}

The rest of the paper is organised as follows. Section~\ref{sec:related} reviews FFIA, audio-visual
fusion, hypercomplex networks and efficient models. Section~\ref{sec:method} presents HyperAVNet,
Section~\ref{sec:setup} the experimental set-up, and Section~\ref{sec:results} the results and
ablations.

%  Contribution bullets that depend on results carry \todo notes: keep, rephrase or drop
%  them once Tables in sections/results.tex are filled.
% ============================================================================
\section{Introduction}
\label{sec:intro}

Feed is one of the largest operating costs in aquaculture, and how it is dispensed directly affects
productivity, water quality and fish welfare~\cite{atoum2014feeding,yang2021smartfish}. Most feeding
machines still run on fixed schedules or on the judgement of an operator, so fish are underfed when
appetite is high and overfed when it is low. The uneaten feed is wasted and degrades the water.
\emph{Fish feeding intensity assessment} (FFIA) estimates the appetite of a group of fish during feeding,
usually as an ordinal level from \textsc{None} to \textsc{Strong}. It is the sensing step needed to close
the loop between feeding machines and the fish~\cite{an2021behaviour,zhou2019ffia}.

Most FFIA systems are vision-based. They classify single images~\cite{zhou2019ffia,yang2021effnetb2,zhang2023msif}
or exploit motion through optical flow and 3-D convolutions~\cite{ubina2021ffia}. Vision captures
aggregation and surface agitation directly, but it suffers from low light, reflections, turbidity
and occlusion, and video models are expensive to run on site. Acoustic FFIA uses the sounds of feeding,
such as chewing, tail slapping and splashing, picked up by a hydrophone~\cite{cui2022affia3k}. It is
unaffected by lighting and cheap to process, but it is sensitive to the bubble and pump noise that is
ever-present in recirculating systems. The two modalities fail under different conditions, which
makes them natural partners. The release of AV-FFIA~\cite{cui2024uffia}, 27,000 synchronised
audio-visual clips of \textit{Oplegnathus punctatus} labelled into four intensities, made this
combination measurable. U-FFIA showed that fusion beats either modality, especially in noise, and
that a unified model can lower the cost~\cite{cui2024uffia}.

Accuracy on AV-FFIA has since risen quickly, with later works reporting 95--98\%
\cite{acafnet2025,dynagate2026,hmmfnet2026,hmkd2026,ldmf2026}. Three problems remain open, however.
\emph{First, efficiency.} The fusion models of~\cite{cui2024uffia} cost 15--69\,GMACs and 21--48\,M
parameters per two-second clip. Recent compact models help, but parameter count, memory and the
cost of transmitting models to edge devices on farms still matter.
\emph{Second, how and where to fuse.} Fusion is usually done once, by concatenation, cross-attention
or attention bottlenecks, at a fixed depth. On AV-FFIA these choices behave very differently: the
attention bottleneck (MBT~\cite{nagrani2021mbt}) reached only 72.4\% against 91.7\% for cross-attention
in~\cite{cui2024uffia}, and there has been no systematic study of the depth at which audio and video
should interact.
\emph{Third, evaluation.} Published numbers rest on incompatible protocols: the official random split,
5-fold cross-validation, and filtered subsets of the data. Moreover, every split is drawn at clip level,
although the clips are consecutive two-second segments of one-minute recordings. Near-duplicate
neighbours can therefore appear in both training and test, which may inflate accuracy.

Hypercomplex neural networks offer a principled answer to the first two problems.
A layer over a hypercomplex algebra (complex numbers, quaternions, octonions, sedenions) shares one set
of weight blocks across all components through the algebra's multiplication table. This divides the
parameter count by the algebra's dimension~$n$ while still mixing every component with every other
\cite{parcollet2019qrnn,zhang2021phm,grassucci2022phnn}. In spatiotemporal tasks such networks match
or exceed their real-valued counterparts at a fraction of the parameters
\cite{bojesomo2024hypercomplex,bojesomo2023octonion,bojesomo2020sedenion}. The same property suggests a
new view of fusion. If the audio view, the video view and their interaction are taken as the
\emph{components} of one hypercomplex number, a single Hamilton product mixes all of them through
shared weights. Cross-modal interaction then becomes a property of the algebra rather than something
a dense layer must learn from limited data.

We build on this view and propose \textbf{HyperAVNet}, a hypercomplex audio-visual network for FFIA. Its
contributions are as follows.
\begin{itemize}
  \item A \textbf{fully hypercomplex audio-visual architecture}: Hamilton-product weight sharing in the
  convolutional backbones, attention, fusion and heads. It has 2.14\,M parameters, $3.7\times$ fewer than
  the identical real-valued network, at 1.18\,GMACs per clip; that is $10\times$ fewer parameters and
  $13\times$ fewer MACs than U-FFIA.
  \item \textbf{Hypercomplex component fusion (HCF)}: modality-specific views and their multiplicative
  agreement are encoded as the components of a quaternion (or octonion) and mixed by one Hamilton
  product. Its parameter cost does not depend on the number of interaction terms.
  \item A \textbf{deep-to-early fusion study with zero-initialised cross-modal feedback}, which shows where
  audio-visual interaction helps on AV-FFIA. \todo{one-line finding, e.g. ``fusing from the second stage
  onwards is best''}
  \item \textbf{Multi-level distillation} at the logit, embedding and relation levels from pretrained
  audio and video teachers, together with a \textbf{teacher ablation} over ten backbone pairs that
  tests whether a stronger teacher yields a stronger student. \todo{one-line finding}
  \item A \textbf{reproducible evaluation} under three protocols (the official split, 5-fold
  cross-validation, and a new \emph{session-disjoint} split), with a noise and missing-modality
  robustness study. \todo{State the headline: accuracy relative to the strongest comparable result,
  and the gap between clip-level and session-disjoint accuracy.}
\end{itemize}

The rest of the paper is organised as follows. Section~\ref{sec:related} reviews FFIA, audio-visual
fusion, hypercomplex networks and efficient models. Section~\ref{sec:method} presents HyperAVNet,
Section~\ref{sec:setup} the experimental set-up, and Section~\ref{sec:results} the results and
ablations.
